\section{One task, complete resources, and attained error alternatives}\label{sec:resources}
\subsection{Task-compatible causal morphisms}
An experiment morphism consists of a parameter-independent causal simulator, a legal strategy lift, report and task-readout transformations, a clock map, and a declared finite internal state. Each virtual report must be emitted before the target decision using it. Raw acquisitions and physical time are mapped as cumulative nonnegative costs. A report transformation without a strategy lift or without a compatible physical audit is not such a morphism. The simulator may not use the later environment-owned audit to generate earlier reports.

\begin{proposition}[Composition in a common bounded task]\label{prop:composition}
Suppose a target observer uses at most $M$ persistent states and a task-compatible simulator at most $S$ internal states. Suppose the complete joint laws of the stopped transcript, the observer's fresh randomization, its committed answer and the protected scored audit differ by total variation at most $\varepsilon$. The composed observer uses at most $MS$ states and, for loss in $[0,L]$, its risk differs by at most $L\varepsilon$. For a chain of simulators, errors add (capped at one), state counts multiply, and the declared call/clock maps compose. A reverse risk inequality requires a reverse task-compatible simulator.
\end{proposition}
\begin{proof}
The runtime state is the pair of observer and simulator labels. The legal strategy lift and report deadlines make its transitions causal. The conditional distribution of each fresh randomization can be applied to a common coupling without increasing total variation. Thus a coupling fails for the complete scored joint law with probability at most $\varepsilon$, and the bounded loss difference is at most $L\varepsilon$. The union bound and induction give the chain statement. No inversion of a one-way simulation is used. A statistic in $[-1,1]$ instead has expectation discrepancy bounded by $2\varepsilon$ under the convention $\TV(P,Q)=\sup_E|P(E)-Q(E)|$.
\end{proof}
Persistent bits cost $\lceil\log_2(MS)\rceil$, not the simulator's program description or its temporary numerical workspace. Internal counters belong to $S$ or $M$, never to an uncharged clock map. This finite-horizon proposition is not a stopping-time, infinite-horizon, pressure or large-deviation equivalence without further hypotheses.

\subsection{Certified finite-precision execution}
\begin{proposition}[Precision with a boundary certificate]\label{prop:precision}
Fix a finite ideal program from the proof of Theorem~\ref{thm:main}. For every phase $t$ and every positive-probability incoming label, suppose the ideal actual pre-quantization row law assigns mass at most $\omega_t(h)$ to the $h$-neighborhood of its cell boundaries. Suppose a digital implementation, while its old label agrees with the ideal one, computes that cell decision with coordinate error at most $h_t$ outside an event of unconditional probability at most $\xi_t$. Suppose it approximates each final readout within $v$. If the primitive kernels or calibration law are additionally coupled with the ideal complete experiment with error $\zeta$, then
\begin{equation}\label{eq:precision}
 |R_{\rm digital}-R_{\rm ideal}|
 \le L\min\{1,\zeta+\sum_t[\omega_t(h_t)+\xi_t]\}+C v
\end{equation}
for bounded squared tasks and a fixed bounded response set. No extra persistent states are required if the compiled label set is unchanged. Numerical workspace, report precision and stored constants must nevertheless be paid separately.
\end{proposition}
\begin{proof}
Use the same reports and fresh randomness until the first label mismatch. Outside the primitive coupling failure and numerical exceptional events, that first mismatch can occur only when the ideal pre-quantization value lies within $h_t$ of a boundary. The unconditional mass of that event is bounded by the row certificate after summing over old labels. Restricting to histories before the first mismatch only decreases this unconditional probability; no density bound conditioned on that restriction is assumed. The union bound gives the displayed minimum. If labels never disagree, the bounded squared loss is Lipschitz in the final response and changes by at most $Cv$. This proves~\eqref{eq:precision} without claiming TV proximity between a continuous report and its quantized numerical representation.
\end{proof}
For a Euclidean grid in a bounded $d$-dimensional coordinate body and actual row densities bounded by $A$, a mesh with $O(m^{1/d})$ hyperplanes has $\omega_t(h)\le C A m^{1/d}h$, capped at one. A Voronoi implementation can instead bound each of its finitely many bisector slabs, giving the conservative certificate $C A m^2h$. Boundary tie rules are fixed Borel rules. In singular components one must supply the corresponding component-wise boundary certificate or choose chart grids; ambient Lebesgue density is not available for free.

The raw Gaussian and positive-effect models with computable, quantitatively separated parameters admit explicit certified arithmetic. For Gaussian reports, truncate computation outside a radius $R$ only as a \emph{charged failure event}: a Gaussian tail bound on the known finite mixture makes its total probability at most a chosen $\xi$, uniformly over the incoming labels. For example $R=C(1+\sqrt{\log(n/\xi)})$ suffices with parameter-dependent constants. The likelihood ratios are evaluated in log coordinates on that bounded event. Positive spectral floors in the quantum model give certified inverses and square roots by polynomial or interval iteration. Rational codepoints can be projected into the declared positive compact domain with a certified tolerance. These operations approximate the chosen program, not an unobservable true parameter.

A deliberately conservative accounting for codebook implementations with $k$ coordinates, $M$ total labels and $p$-bit constants is: $O(kMp)$ read-only codepoint bits, the separately specified primitive parameter description, a transition/phase description of $O(M\log M)$ bits, and a final response table included in the same $O(kMp)$ term. Naive nearest-center evaluation uses $O(Mk)$ arithmetic operations per report in addition to primitive evaluation; $O(k^2p+\log M)$ working bits suffice for serial distance and matrix calculations at a declared internal accuracy, with any larger interval-iteration workspace stated explicitly. Counts refer to shared read-only constants rather than one copy per traversed edge. The report-precision oracle and its time are part of the implementation interface. Calibrated-cell construction under (C) additionally requires certified forward integrals for its representatives; mere Borel existence does not supply those computations. No matched program-length, workspace or arithmetic-time lower is inferred from the label theorem.

\subsection{A matching calibration--deficiency law with a counted retained bit}
A positive upper error budget by itself cannot supply a minimax lower. We next give an attained physical alternative and count the memory needed to retain its useful report. Fix any core experiment of Theorem~\ref{thm:main} in dimension $d$, with $n\ge1$ core report calls. Independently prepare a fair bit $U\in\{-1,1\}$ and an unknown fixed calibration offset $\theta\in\{-\delta,\delta\}$. The initial paid action \textsc{Bit} reports $U$ with probability $1-a$ and the erasure symbol with probability $a$, where $0\le a\le1/2$. It must be followed by exactly $n$ passive core \textsc{Filter} calls. The answer is committed, and then one joint audit returns $(Z_n,U,\theta)$. The loss is the sum of the three squared coordinate losses. The true sign of $\theta$ is unavailable to the program. The preparation, bit call, core calls and final audit are
\begin{equation}\label{eq:physical-N}
 N=n+3
\end{equation}
physical calls. The initial action and the later \textsc{Filter} actions must be selected by the machine's labels; no action-availability signal supplies a phase. Programs issuing a different prescribed action word are not legal for this experiment.

Compare this experiment to the target whose initial bit report is exact, with the same protected audit and core experiment. The simulator must emit its virtual bit before any core reports. Put $\varepsilon=a/2$.

\begin{theorem}[Attained joint law and full memory charge]\label{thm:joint}
The uniform causal deficiency of the erased-bit interface to the exact-bit target, with its audit protected and its report deadline fixed, is exactly $\varepsilon$. For any fixed bounded range of $\delta$, the joint minimax checkpoint risk satisfies
\begin{equation}\label{eq:joint-cp}
 \calR^{\rm joint}_{\rm cp}(n,M,\delta,\varepsilon)-B_n
 \asymp M^{-2/d}+\delta^2+\varepsilon,
 \qquad M\ge1.
\end{equation}
The legal autonomous class is nonempty exactly when $M\ge n+2=N-1$. On that class, set
\[
 Q(n,M)=\max\{1,M-2n-1\}.
\]
Then
\begin{equation}\label{eq:joint-aut}
 \calR^{\rm joint}_{\rm aut}(n,M,\delta,\varepsilon)-B_n
 \asymp Q(n,M)^{-2/d}+\delta^2+\varepsilon.
\end{equation}
Constants depend only on the core certificate, the task bound and the fixed offset range; $a$ itself is retained in the risk. Neither the true bit on erasure nor the calibration sign is passed to the observer. In terms of physical acquisitions, substitute $n=N-3$ in~\eqref{eq:joint-aut}.
\end{theorem}
\begin{proof}
On a nonerasure, transmit the bit; on erasure, emit a fresh fair bit. This stateless simulator has disagreement probability $a/2$ with the physical audit bit, and otherwise matches the target joint law, including all later core observations. Its total variation error is $a/2$. Conversely, conditional on erasure, every admissible simulator's pre-core information is independent of $U$, even with arbitrary additional classical memory. It guesses incorrectly with probability $1/2$. The event that its virtual bit disagrees with the later protected physical bit has probability at least $a/2$ and probability zero under the exact target. This proves the exact deficiency. Allowing the simulator to replace that audit would be a different comparison.

For either observer class, averaging the squared offset loss over the two signs gives at least $\delta^2$; the response zero attains this. Even the complete available history has bit error $a$ on erasure and zero otherwise. The independent core's all-label last-row converse remains valid with the bit and observer randomizations included in the old label. It yields $cM^{-2/d}$ at a checkpoint. Summing the three lower bounds under the common uniform sign prior, and then using minimax domination, proves the lower in~\eqref{eq:joint-cp}. For the upper, impute an independent fair bit on erasure and keep the resulting sign $b$. The fixed response $(1-a)b$ has bit risk $2a-a^2\le2a$. Use two copies of a core checkpoint code with at most $\lfloor M/2\rfloor$ labels. For $M\ge2$ this gives the upper, and for $M=1$ a constant response gives the same order after enlarging constants.

For an autonomous machine, the initial \textsc{Bit} label cannot recur at a \textsc{Filter} phase because its action row is fixed. All subsequent core-phase report laws have common support, so Lemma~\ref{lem:clock} separates the positive-probability label sets at the phases after the bit and after each core call, including halt. Thus $n+2$ labels are necessary and one label per phase suffices for legality.

There are now two possibilities. If any of these post-bit phase sets, including halt, contains only one positive-probability label, all retained information about $U$ has been lost at that phase. The physical core and its future reports are independent of $U$, and there is no external decoder transcript or persistent seed. The bit error at the final response is therefore at least one. If every nonterminal post-bit phase contains at least two labels, they use at least $2n$ labels, besides the initial label. Hence the final phase has at most $M-2n-1$ labels, and the core last-row converse yields $cQ(n,M)^{-2/d}$. If its final phase instead has one label, the first case applies. In either case the joint risk above $B_n$ is bounded below by a constant times $Q^{-2/d}+a$, plus $\delta^2$: in the first case $Q^{-2/d}\le1$ and $a\le1/2$; in the second use the core lower and the complete-history bit lower. This also covers randomized encoders, since all phase sets and independence statements are under their actual laws.

For the upper when $M\le2n+2$, ignore the bit and use a legal constant-output core program, giving bounded excess; here $Q=1$. For $M\ge2n+3$, reserve one initial label, two imputed-bit labels after the first call, and two copies of every later core-phase label. Choose widths $k_t$ with
\[
 1+2+2\sum_{t=1}^nk_t\le M.
\]
The core construction may therefore use total budget
$M_0=1+\lfloor(M-3)/2\rfloor\ge n+1$. Its excess is at most $C(M_0-n)^{-2/d}$, and $M_0-n$ is comparable to $Q$ by fixed constants. Keep the sign tag throughout and output $(1-a)b$, and output zero for the offset. The resulting excess is at most $CQ^{-2/d}+2a+\delta^2$. Every tag, phase and halt label has been charged. This proves~\eqref{eq:joint-aut}.
\end{proof}
The earlier articles already use the elementary erasure and hidden-offset devices. The new assertion here is the two-sided retained-bit clock threshold and its combination with the broader Bayesian filter class. Formula~\eqref{eq:joint-aut} is an attained same-task law, not a universal lower bound for every calibration or simulator error of those sizes. The exact count $N$ is a prescribed exposure protocol for a changing physical state, not an unrestricted at-most-$N$ static experimental-design problem.
